\section{Witnessing quantum entanglement (9/18)}

\newcommand{\epr}{\mathrm{EPR}}
\newcommand{\ghz}{\mathrm{GHZ}}

\referto{\tong~16.1.1-16.1.4, \scully~Chapter 18}



\subsection{The EPR ``paradox''}

Einstein, Podolsky and Rosen (EPR) investigated the singlet Bell state of 2 qubits
\begin{align}
    \dket\epr=\frac1{\sqrt2}(\dket{01}-\dket{10})
\end{align}
which was also known as the EPR state. Assuming that Alice has qubit 1 and Bob has qubit 2, they observed that a measurement of Pauli $\hat Z_1$ operator by Alice (say, with outcome $q=\pm1$) will instantaneously collapse the EPR state into a product state $\dket{q}_1\otimes\dket{-q}_2$, and immediately determine the measurement outcome of Pauli $\hat Z_2$ operator by Bob, no matter how far away Alice and Bob are physically separated. 

EPR considered the following 3 claims which all seemed reasonable, from the perspective of classical physics:

\textbf{Locality:} Alice’s choice cannot instantaneously alter the physical reality at Bob’s distant location.

\textbf{Reality criterion:} Perfectly predictable results must reflect preexisting properties.

\textbf{Quantum completeness:} The wavefunction completely describes physical reality.

EPR concluded that these three claims cannot all hold simultaneously. Because they retained locality and the reality criterion, they concluded that the quantum state is incomplete, perhaps omitting hidden variables that specify the missing spin values.

The Einstein--Podolsky--Rosen paradox asks whether quantum mechanics can
be replaced by a theory that is both local and realistic. Consider two
spatially separated observers, Alice and Bob, who receive qubits $1$ and $2$.
Alice may measure either $Q$ or $R$, while Bob may measure either $S$ or $T$.
Every measurement has outcomes in $\{+1,-1\}$:
\begin{equation}
 Q^2=R^2=S^2=T^2=1.
\end{equation}
Measurements associated with different parties are compatible, for example
$[Q,S]=0$.

In a local realistic model, measurement outcomes are assumed to have
pre-existing values and Alice's choice does not affect Bob's value. Thus, for
a hidden-variable assignment,
\begin{equation}
 E(QS)=E(Q)E(S),
\end{equation}
where each individual value $E(Q),E(R),E(S),E(T)$ is $\pm1$. 

\subsection{The CHSH inequality}
Define the CHSH combination
\begin{equation}
 \mathcal{B}=QS+RS+RT-QT.
\end{equation}
For a local realistic assignment,
\begin{align}
 E(\mathcal{B})
 &=\bigl[E(Q)+E(R)\bigr]E(S)
   +\bigl[E(R)-E(Q)\bigr]E(T).
\end{align}
Since $E(R)=E(Q)$ or $E(R)=-E(Q)$, precisely one bracket vanishes and the
other equals $\pm2$. Consequently,
\begin{equation}
 \boxed{|E(QS+RS+RT-QT)|\leq2.}
\end{equation}
This is the Clauser--Horne--Shimony--Holt (CHSH) inequality.

\subsubsection{Quantum-mechanical violation}
Let the observables be spin measurements along unit vectors $\bm m_{1,2}$ and $\bm n_{1,2}$:
\begin{equation}
 Q=\bm m_1\cdot\bm\sigma_1,
 \quad R=\bm n_1\cdot\bm\sigma_1,
 \quad S=\bm m_2\cdot\bm\sigma_2,
 \quad T=\bm n_2\cdot\bm\sigma_2.
\end{equation}
The two qubits are prepared in the singlet Bell state
\begin{equation}
 \dket{\Psi^-}=\frac{1}{\sqrt2}\left(\dket{01}-\dket{10}\right).
\end{equation}
For this state,
\begin{equation}
 \dbra{\Psi^-}(\bm a\cdot\bm\sigma_1)
 (\bm b\cdot\bm\sigma_2)\dket{\Psi^-}=-\bm a\cdot\bm b.
\end{equation}
Hence
\begin{align}
 \langle\mathcal{B}\rangle
 &=-\bm m_1\cdot\bm m_2-\bm n_1\cdot\bm m_2
   -\bm n_1\cdot\bm n_2+\bm m_1\cdot\bm n_2\\
 &=-(\bm m_1+\bm n_1)\cdot\bm m_2
   +(\bm m_1-\bm n_1)\cdot\bm n_2.
\end{align}
Choose coplanar directions with neighboring axes separated by $45^\circ$.
Then
\begin{equation}
 |\langle\mathcal{B}\rangle|=2\sqrt2>2.
\end{equation}
The value $2\sqrt2$ is the maximal quantum value, known as the Tsirelson
bound.

\subsubsection{Experimental test with photon polarization}
In an optical realization, a source produces pairs of entangled photons. One
photon is sent to polarimeter $A$ and the other to polarimeter $B$. Each
polarimeter selects a measurement orientation, and a coincidence monitor
records joint detection events. With
\begin{equation}
 \dket{V}=\text{vertical polarization},
 \qquad \dket{H}=\text{horizontal polarization},
\end{equation}
polarization measurements implement the two-outcome observables required for
a Bell test.

\subsection{The GHZ argument}

Consider the three-qubit GHZ state
\begin{equation}
 \dket{\ghz}=\frac{1}{\sqrt2}\left(\dket{000}-\dket{111}\right).
\end{equation}
The state obeys
\begin{align}
 X_A Y_B Y_C\dket{\ghz}&=+\dket{\ghz},\\
 Y_A X_B Y_C\dket{\ghz}&=+\dket{\ghz},\\
 Y_A Y_B X_C\dket{\ghz}&=+\dket{\ghz},
\end{align}
and
\begin{equation}
 X_A X_B X_C\dket{\ghz}=-\dket{\ghz}.
\end{equation}

Suppose a local hidden-variable theory assigns predetermined values
$s_j^{(x)},s_j^{(y)}\in\{+1,-1\}$ to the $X$ and $Y$ measurements on qubit
$j=A,B,C$. The first three GHZ correlations require
\begin{align}
 s_A^{(x)}s_B^{(y)}s_C^{(y)}&=+1,\\
 s_A^{(y)}s_B^{(x)}s_C^{(y)}&=+1,\\
 s_A^{(y)}s_B^{(y)}s_C^{(x)}&=+1.
\end{align}
Multiplying these equations gives
\begin{equation}
 s_A^{(x)}s_B^{(x)}s_C^{(x)}=+1.
\end{equation}
Quantum mechanics instead predicts a product of $-1$. This is a direct
contradiction, with no statistical inequality required.

\subsection{Bell's Original Inequality}

Return to the singlet state
\begin{equation}
 \dket{\epr}=\frac{1}{\sqrt2}\left(\dket{01}-\dket{10}\right).
\end{equation}
For a unit vector $\bm n$,
\begin{equation}
 \bm n\cdot\bm\sigma=n_x\sigma_x+n_y\sigma_y+n_z\sigma_z,
 \qquad |\bm n|=1,
\end{equation}
and
\begin{equation}
 \bm n\cdot\bm\sigma\,\dket{\pm_{\bm n}}
 =\pm\dket{\pm_{\bm n}}.
\end{equation}
In any common spin basis,
\begin{equation}
 \dket{\epr}=\frac{1}{\sqrt2}\left(
 \dket{+_{\bm n},-_{\bm n}}-\dket{-_{\bm n},+_{\bm n}}
 \right).
\end{equation}

Consider measurement axes $\bm n_a$, $\bm n_b$, and $\bm n_c$. Let $P_{ab}$
denote the probability that the first qubit has outcome $+1$ along
$\bm n_a$ while the second has outcome $+1$ along $\bm n_b$. Writing
$p(s_as_bs_c)$ for a predetermined hidden-variable assignment,
\begin{align}
 P_{ab}&=p(+,-,+)+p(+,-,-),\\
 P_{bc}&=p(+,+,-)+p(-,+,-),\\
 P_{ac}&=p(+,-,-)+p(-,+,-).
\end{align}
Nonnegativity of probabilities implies a three-axis probability inequality
\begin{equation}
 \boxed{P_{ab}+P_{bc}\geq P_{ac}.}
\end{equation}

This inequality is indeeded violated by quantum probabilities. For coplanar directions, parameterize the spin eigenstates by
\begin{equation}
 \dket{+_{\theta}}=
 \begin{pmatrix}\cos(\theta/2)\\ \sin(\theta/2)\end{pmatrix},
 \qquad
 \dket{-_{\theta}}=
 \begin{pmatrix}-\sin(\theta/2)\\ \cos(\theta/2)\end{pmatrix}.
\end{equation}
For the singlet,
\begin{equation}
 P_{ij}=\frac12\sin^2\!\left(\frac{\theta_i-\theta_j}{2}\right).
\end{equation}
Choose
\begin{equation}
 \theta_a=0,\qquad \theta_b=\theta,\qquad \theta_c=2\theta.
\end{equation}
Then
\begin{equation}
 P_{ab}=P_{bc}=\frac12\sin^2\!\left(\frac\theta2\right),
 \qquad P_{ac}=\frac12\sin^2\theta.
\end{equation}
Bell's inequality would require
\begin{equation}
 \sin^2\!\left(\frac\theta2\right)\geq\frac12\sin^2\theta.
\end{equation}
For $\theta=60^\circ$,
\begin{equation}
 \frac14\not\geq\frac38,
\end{equation}
which contradicts the local hidden-variable prediction.

\subsection{Summary}

Below is a summary of these entanglement witnesses:
\begin{itemize}
 \item Local realism leads to the CHSH bound
       $|\langle\mathcal B\rangle|\leq2$.
 \item The singlet state can reach $2\sqrt2$, violating the CHSH inequality.
 \item GHZ correlations give a deterministic contradiction with predetermined
       local outcomes.
 \item Bell's original probability inequality is also violated by singlet
       correlations for suitable measurement directions.
\end{itemize}

